% Options for packages loaded elsewhere
\PassOptionsToPackage{unicode}{hyperref}
\PassOptionsToPackage{hyphens}{url}
%
\documentclass[
  ignorenonframetext,
]{beamer}
\usepackage{pgfpages}
\setbeamertemplate{caption}[numbered]
\setbeamertemplate{caption label separator}{: }
\setbeamercolor{caption name}{fg=normal text.fg}
\beamertemplatenavigationsymbolsempty
% Prevent slide breaks in the middle of a paragraph
\widowpenalties 1 10000
\raggedbottom
\setbeamertemplate{part page}{
  \centering
  \begin{beamercolorbox}[sep=16pt,center]{part title}
    \usebeamerfont{part title}\insertpart\par
  \end{beamercolorbox}
}
\setbeamertemplate{section page}{
  \centering
  \begin{beamercolorbox}[sep=12pt,center]{section title}
    \usebeamerfont{section title}\insertsection\par
  \end{beamercolorbox}
}
\setbeamertemplate{subsection page}{
  \centering
  \begin{beamercolorbox}[sep=8pt,center]{subsection title}
    \usebeamerfont{subsection title}\insertsubsection\par
  \end{beamercolorbox}
}
\AtBeginPart{
  \frame{\partpage}
}
\AtBeginSection{
  \ifbibliography
  \else
    \frame{\sectionpage}
  \fi
}
\AtBeginSubsection{
  \frame{\subsectionpage}
}
\usepackage{amsmath,amssymb}
\usepackage{iftex}
\ifPDFTeX
  \usepackage[T1]{fontenc}
  \usepackage[utf8]{inputenc}
  \usepackage{textcomp} % provide euro and other symbols
\else % if luatex or xetex
  \usepackage{unicode-math} % this also loads fontspec
  \defaultfontfeatures{Scale=MatchLowercase}
  \defaultfontfeatures[\rmfamily]{Ligatures=TeX,Scale=1}
\fi
\usepackage{lmodern}
\ifPDFTeX\else
  % xetex/luatex font selection
\fi
% Use upquote if available, for straight quotes in verbatim environments
\IfFileExists{upquote.sty}{\usepackage{upquote}}{}
\IfFileExists{microtype.sty}{% use microtype if available
  \usepackage[]{microtype}
  \UseMicrotypeSet[protrusion]{basicmath} % disable protrusion for tt fonts
}{}
\makeatletter
\@ifundefined{KOMAClassName}{% if non-KOMA class
  \IfFileExists{parskip.sty}{%
    \usepackage{parskip}
  }{% else
    \setlength{\parindent}{0pt}
    \setlength{\parskip}{6pt plus 2pt minus 1pt}}
}{% if KOMA class
  \KOMAoptions{parskip=half}}
\makeatother
\usepackage{xcolor}
\newif\ifbibliography
\usepackage{longtable,booktabs,array}
\usepackage{calc} % for calculating minipage widths
\usepackage{caption}
% Make caption package work with longtable
\makeatletter
\def\fnum@table{\tablename~\thetable}
\makeatother
\usepackage{graphicx}
\makeatletter
\newsavebox\pandoc@box
\newcommand*\pandocbounded[1]{% scales image to fit in text height/width
  \sbox\pandoc@box{#1}%
  \Gscale@div\@tempa{\textheight}{\dimexpr\ht\pandoc@box+\dp\pandoc@box\relax}%
  \Gscale@div\@tempb{\linewidth}{\wd\pandoc@box}%
  \ifdim\@tempb\p@<\@tempa\p@\let\@tempa\@tempb\fi% select the smaller of both
  \ifdim\@tempa\p@<\p@\scalebox{\@tempa}{\usebox\pandoc@box}%
  \else\usebox{\pandoc@box}%
  \fi%
}
% Set default figure placement to htbp
\def\fps@figure{htbp}
\makeatother
\setlength{\emergencystretch}{3em} % prevent overfull lines
\providecommand{\tightlist}{%
  \setlength{\itemsep}{0pt}\setlength{\parskip}{0pt}}
\setcounter{secnumdepth}{-\maxdimen} % remove section numbering
\usepackage{booktabs}
\usepackage{longtable}
\usepackage{array}
\usepackage{multirow}
\usepackage{wrapfig}
\usepackage{float}
\usepackage{colortbl}
\usepackage{pdflscape}
\usepackage{tabu}
\usepackage{threeparttable}
\usepackage{threeparttablex}
\usepackage[normalem]{ulem}
\usepackage{makecell}
\usepackage{xcolor}
\usepackage{multicol}
\usepackage{ulem}
\usepackage{hhline}
\newlength\Oldarrayrulewidth
\newlength\Oldtabcolsep
\usepackage{hyperref}
\usepackage{bookmark}
\IfFileExists{xurl.sty}{\usepackage{xurl}}{} % add URL line breaks if available
\urlstyle{same}
\hypersetup{
  pdftitle={NSW Road Crash Severity Prediction},
  pdfauthor={Sanskar \textbar{} Robert \textbar{} Haipeng \textbar{} Chen \textbar{} Yihan},
  hidelinks,
  pdfcreator={LaTeX via pandoc}}

\title{NSW Road Crash Severity Prediction}
\subtitle{STAT5003 - Group 07, Workshop 12}
\author{Sanskar \textbar{} Robert \textbar{} Haipeng \textbar{} Chen
\textbar{} Yihan}
\date{21 May 2026}

\begin{document}
\frame{\titlepage}

\begin{frame}
\end{frame}

\begin{frame}{Research Question \& Motivation}
\phantomsection\label{research-question-motivation}
\textbf{Research Question:}

\begin{quote}
Given environmental conditions (weather, lighting, road surface, speed
limit) and traffic unit characteristics (TU type, manoeuvre, number of
units), can we predict NSW road crash severity \textbf{(Fatal / Injury /
Non-casualty)}? A three-class classification problem.
\end{quote}

\textbf{Why this matters:}

\begin{itemize}
\tightlist
\item
  NSW recorded \textbf{340 road fatalities in 2024}; over two-thirds in
  regional areas \emph{(TfNSW, 2025)}
\item
  Fatal crashes are only \textbf{\textasciitilde1.4\%} of all records -
  a severe class imbalance that naive models simply ignore
\item
  Environmental and traffic conditions are knowable \emph{at the time of
  the crash}, making real-time prediction feasible
\item
  A robust model provides actionable intelligence across multiple
  decision-making contexts
\end{itemize}

\textbf{Stakeholder value:}

\begin{longtable}[]{@{}
  >{\raggedright\arraybackslash}p{(\linewidth - 2\tabcolsep) * \real{0.5000}}
  >{\raggedright\arraybackslash}p{(\linewidth - 2\tabcolsep) * \real{0.5000}}@{}}
\toprule\noalign{}
\begin{minipage}[b]{\linewidth}\raggedright
Stakeholder
\end{minipage} & \begin{minipage}[b]{\linewidth}\raggedright
Application
\end{minipage} \\
\midrule\noalign{}
\endhead
\textbf{Traffic Authorities} & Target speed limits and road design in
high-risk corridors \\
\textbf{Emergency Responders} & Dispatch resources proportional to
predicted severity before on-scene confirmation \\
\textbf{Policymakers \& Insurers} & Evidence-based road safety policy
and risk-adjusted pricing \\
\bottomrule\noalign{}
\end{longtable}

The core challenge: Fatal crashes are rare (\textasciitilde1.4\%) but
highest-consequence. Standard accuracy metrics fail here - we need
models and evaluation that explicitly prioritise the minority class.
\end{frame}

\begin{frame}{EDA \& Data Understanding: Dataset Overview}
\phantomsection\label{eda-data-understanding-dataset-overview}
\global\setlength{\Oldarrayrulewidth}{\arrayrulewidth}

\global\setlength{\Oldtabcolsep}{\tabcolsep}

\setlength{\tabcolsep}{2pt}

\renewcommand*{\arraystretch}{1.5}



\providecommand{\ascline}[3]{\noalign{\global\arrayrulewidth #1}\arrayrulecolor[HTML]{#2}\cline{#3}}

\begin{longtable}[c]{cc}



\hhline{>{\arrayrulecolor[HTML]{666666}\global\arrayrulewidth=1.5pt}->{\arrayrulecolor[HTML]{666666}\global\arrayrulewidth=1.5pt}-}

\multicolumn{1}{>{}c}{\textcolor[HTML]{000000}{\fontsize{18}{18}\selectfont{\textbf{Item}}}} & \multicolumn{1}{>{}c}{\textcolor[HTML]{000000}{\fontsize{18}{18}\selectfont{\textbf{Detail}}}} \\

\noalign{\global\arrayrulewidth 0pt}\arrayrulecolor[HTML]{000000}

\hhline{>{\arrayrulecolor[HTML]{666666}\global\arrayrulewidth=1.5pt}->{\arrayrulecolor[HTML]{666666}\global\arrayrulewidth=1.5pt}-}\endfirsthead 

\hhline{>{\arrayrulecolor[HTML]{666666}\global\arrayrulewidth=1.5pt}->{\arrayrulecolor[HTML]{666666}\global\arrayrulewidth=1.5pt}-}

\multicolumn{1}{>{}c}{\textcolor[HTML]{000000}{\fontsize{18}{18}\selectfont{\textbf{Item}}}} & \multicolumn{1}{>{}c}{\textcolor[HTML]{000000}{\fontsize{18}{18}\selectfont{\textbf{Detail}}}} \\

\noalign{\global\arrayrulewidth 0pt}\arrayrulecolor[HTML]{000000}

\hhline{>{\arrayrulecolor[HTML]{666666}\global\arrayrulewidth=1.5pt}->{\arrayrulecolor[HTML]{666666}\global\arrayrulewidth=1.5pt}-}\endhead



\multicolumn{1}{>{}c}{\textcolor[HTML]{000000}{\fontsize{16}{16}\selectfont{\textbf{Time\ Period}}}} & \multicolumn{1}{>{}c}{\textcolor[HTML]{000000}{\fontsize{16}{16}\selectfont{2020\ -\ 2024}}} \\

\noalign{\global\arrayrulewidth 0pt}\arrayrulecolor[HTML]{000000}

\hhline{>{\arrayrulecolor[HTML]{666666}\global\arrayrulewidth=0.75pt}->{\arrayrulecolor[HTML]{666666}\global\arrayrulewidth=0.75pt}-}



\multicolumn{1}{>{\cellcolor[HTML]{F5F5F5}}c}{\textcolor[HTML]{000000}{\fontsize{16}{16}\selectfont{\textbf{Data\ Sources}}}} & \multicolumn{1}{>{\cellcolor[HTML]{F5F5F5}}c}{\textcolor[HTML]{000000}{\fontsize{16}{16}\selectfont{CRASH.xlsx\ +\ TRAFFIC\ UNIT.xlsx\ (TfNSW\ via\ data.gov.au)}}} \\

\noalign{\global\arrayrulewidth 0pt}\arrayrulecolor[HTML]{000000}

\hhline{>{\arrayrulecolor[HTML]{666666}\global\arrayrulewidth=0.75pt}->{\arrayrulecolor[HTML]{666666}\global\arrayrulewidth=0.75pt}-}



\multicolumn{1}{>{}c}{\textcolor[HTML]{000000}{\fontsize{16}{16}\selectfont{\textbf{Merge\ Key}}}} & \multicolumn{1}{>{}c}{\textcolor[HTML]{000000}{\fontsize{16}{16}\selectfont{Crash\ ID\ (left\ join)}}} \\

\noalign{\global\arrayrulewidth 0pt}\arrayrulecolor[HTML]{000000}

\hhline{>{\arrayrulecolor[HTML]{666666}\global\arrayrulewidth=0.75pt}->{\arrayrulecolor[HTML]{666666}\global\arrayrulewidth=0.75pt}-}



\multicolumn{1}{>{\cellcolor[HTML]{F5F5F5}}c}{\textcolor[HTML]{000000}{\fontsize{16}{16}\selectfont{\textbf{Total\ Records}}}} & \multicolumn{1}{>{\cellcolor[HTML]{F5F5F5}}c}{\textcolor[HTML]{000000}{\fontsize{16}{16}\selectfont{170,962}}} \\

\noalign{\global\arrayrulewidth 0pt}\arrayrulecolor[HTML]{000000}

\hhline{>{\arrayrulecolor[HTML]{666666}\global\arrayrulewidth=0.75pt}->{\arrayrulecolor[HTML]{666666}\global\arrayrulewidth=0.75pt}-}



\multicolumn{1}{>{}c}{\textcolor[HTML]{000000}{\fontsize{16}{16}\selectfont{\textbf{Total\ Variables}}}} & \multicolumn{1}{>{}c}{\textcolor[HTML]{000000}{\fontsize{16}{16}\selectfont{59}}} \\

\noalign{\global\arrayrulewidth 0pt}\arrayrulecolor[HTML]{000000}

\hhline{>{\arrayrulecolor[HTML]{666666}\global\arrayrulewidth=0.75pt}->{\arrayrulecolor[HTML]{666666}\global\arrayrulewidth=0.75pt}-}



\multicolumn{1}{>{\cellcolor[HTML]{F5F5F5}}c}{\textcolor[HTML]{000000}{\fontsize{16}{16}\selectfont{\textbf{Target\ Variable}}}} & \multicolumn{1}{>{\cellcolor[HTML]{F5F5F5}}c}{\textcolor[HTML]{000000}{\fontsize{16}{16}\selectfont{Degree\ of\ Crash\ (3-class\ classification)}}} \\

\noalign{\global\arrayrulewidth 0pt}\arrayrulecolor[HTML]{000000}

\hhline{>{\arrayrulecolor[HTML]{666666}\global\arrayrulewidth=0.75pt}->{\arrayrulecolor[HTML]{666666}\global\arrayrulewidth=0.75pt}-}



\multicolumn{1}{>{}c}{\textcolor[HTML]{000000}{\fontsize{16}{16}\selectfont{\textbf{Fatal}}}} & \multicolumn{1}{>{}c}{\textcolor[HTML]{D62728}{\fontsize{16}{16}\selectfont{2,416\ (1.4\%)}}} \\

\noalign{\global\arrayrulewidth 0pt}\arrayrulecolor[HTML]{000000}

\hhline{>{\arrayrulecolor[HTML]{666666}\global\arrayrulewidth=0.75pt}->{\arrayrulecolor[HTML]{666666}\global\arrayrulewidth=0.75pt}-}



\multicolumn{1}{>{}c}{\textcolor[HTML]{000000}{\fontsize{16}{16}\selectfont{\textbf{Injury}}}} & \multicolumn{1}{>{}c}{\textcolor[HTML]{FF7F0E}{\fontsize{16}{16}\selectfont{114,355\ (66.9\%)}}} \\

\noalign{\global\arrayrulewidth 0pt}\arrayrulecolor[HTML]{000000}

\hhline{>{\arrayrulecolor[HTML]{666666}\global\arrayrulewidth=0.75pt}->{\arrayrulecolor[HTML]{666666}\global\arrayrulewidth=0.75pt}-}



\multicolumn{1}{>{}c}{\textcolor[HTML]{000000}{\fontsize{16}{16}\selectfont{\textbf{Non-casualty}}}} & \multicolumn{1}{>{}c}{\textcolor[HTML]{2CA02C}{\fontsize{16}{16}\selectfont{54,191\ (31.7\%)}}} \\

\noalign{\global\arrayrulewidth 0pt}\arrayrulecolor[HTML]{000000}

\hhline{>{\arrayrulecolor[HTML]{666666}\global\arrayrulewidth=1.5pt}->{\arrayrulecolor[HTML]{666666}\global\arrayrulewidth=1.5pt}-}



\end{longtable}



\arrayrulecolor[HTML]{000000}

\global\setlength{\arrayrulewidth}{\Oldarrayrulewidth}

\global\setlength{\tabcolsep}{\Oldtabcolsep}

\renewcommand*{\arraystretch}{1}

The merged dataset combines crash-level environmental conditions with
traffic unit vehicle data, giving us both where a crash happened and who
was involved. 170,962 records spanning 5 years provide robust signal for
modelling, but the extreme class imbalance (Fatal = 1.4\%) demands
careful handling throughout.
\end{frame}

\begin{frame}{EDA \& Data Understanding: Class Imbalance}
\phantomsection\label{eda-data-understanding-class-imbalance}
\pandocbounded{\includegraphics[keepaspectratio]{STAT5003_W12G07_A2_Presentation_files/figure-beamer/class-imbalance-robert-1.pdf}}

\begin{longtable}[]{@{}
  >{\raggedright\arraybackslash}p{(\linewidth - 2\tabcolsep) * \real{0.6923}}
  >{\raggedright\arraybackslash}p{(\linewidth - 2\tabcolsep) * \real{0.3077}}@{}}
\toprule\noalign{}
\begin{minipage}[b]{\linewidth}\raggedright
Issue / Decision
\end{minipage} & \begin{minipage}[b]{\linewidth}\raggedright
Detail
\end{minipage} \\
\midrule\noalign{}
\endhead
\textbf{Class imbalance} & Fatal only \textasciitilde1.4\% of all
crashes \\
\textbf{Risk of naive models} & Accuracy-maximising models ignore Fatal
entirely \\
\textbf{Imbalance strategy} & Inverse-frequency class weights (all
models) \\
\textbf{SMOTE} & Applied inside training fold only; synthetic Fatal
share raised to \textasciitilde10\% in training \\
\textbf{Primary metrics} & Macro-F1 and Fatal Recall (target
\textgreater= 0.50) \\
\bottomrule\noalign{}
\end{longtable}
\end{frame}

\begin{frame}[fragile]{EDA \& Data Understanding: Speed vs Severity}
\phantomsection\label{eda-data-understanding-speed-vs-severity}
\pandocbounded{\includegraphics[keepaspectratio]{STAT5003_W12G07_A2_Presentation_files/figure-beamer/slide-speed-1.pdf}}

\textbf{Fatal rate at 110 km/h} is \textasciitilde5x higher than at 50
km/h - the single strongest environmental predictor in the dataset.

\textbf{50-60 km/h zones} dominate crash \emph{volume} (urban traffic),
but fatal \emph{proportion} is low - volume and risk are decoupled.

\textbf{Motivates \texttt{speed\_band}} feature: grouping ≤50 / 60-80 /
90-100 / 110 captures the non-linear jump in fatal risk.
\end{frame}

\begin{frame}[fragile]{EDA \& Data Understanding: High-Risk Scenarios}
\phantomsection\label{eda-data-understanding-high-risk-scenarios}
\begingroup\fontsize{13}{15}\selectfont

\begin{longtabu} to \linewidth {>{\raggedright}X>{\raggedright}X>{\raggedright}X>{\raggedleft}X>{}r}
\toprule
\cellcolor[HTML]{2E74B5}{\textcolor{white}{\textbf{Speed Limit}}} & \cellcolor[HTML]{2E74B5}{\textcolor{white}{\textbf{Lighting}}} & \cellcolor[HTML]{2E74B5}{\textcolor{white}{\textbf{Urbanisation}}} & \cellcolor[HTML]{2E74B5}{\textcolor{white}{\textbf{Total Crashes}}} & \cellcolor[HTML]{2E74B5}{\textcolor{white}{\textbf{Fatal \%}}}\\
\midrule
\cellcolor[HTML]{fce4e4}{110 km/h} & \cellcolor[HTML]{fce4e4}{Dusk} & \cellcolor[HTML]{fce4e4}{Country non-urban} & \cellcolor[HTML]{fce4e4}{230} & \textcolor[HTML]{d62728}{\textbf{\cellcolor[HTML]{fce4e4}{7.8\%}}}\\
\cellcolor[HTML]{fce4e4}{100 km/h} & \cellcolor[HTML]{fce4e4}{Darkness} & \cellcolor[HTML]{fce4e4}{Country non-urban} & \cellcolor[HTML]{fce4e4}{2565} & \textcolor[HTML]{d62728}{\textbf{\cellcolor[HTML]{fce4e4}{6.8\%}}}\\
\cellcolor[HTML]{fce4e4}{100 km/h} & \cellcolor[HTML]{fce4e4}{Dusk} & \cellcolor[HTML]{fce4e4}{Country non-urban} & \cellcolor[HTML]{fce4e4}{626} & \textcolor[HTML]{d62728}{\textbf{\cellcolor[HTML]{fce4e4}{6.7\%}}}\\
100 km/h & Daylight & Country non-urban & 7490 & \textcolor[HTML]{d62728}{\textbf{5.5\%}}\\
110 km/h & Darkness & Country non-urban & 1215 & \textcolor[HTML]{d62728}{\textbf{5.4\%}}\\
\addlinespace
80 km/h & Dawn & Country urban & 270 & \textcolor[HTML]{d62728}{\textbf{4.8\%}}\\
70 km/h & Darkness & Country urban & 463 & \textcolor[HTML]{d62728}{\textbf{4.8\%}}\\
80 km/h & Darkness & Country urban & 1536 & \textcolor[HTML]{d62728}{\textbf{4.7\%}}\\
110 km/h & Darkness & Sydney metro. area & 245 & \textcolor[HTML]{d62728}{\textbf{4.5\%}}\\
100 km/h & Dawn & Country non-urban & 455 & \textcolor[HTML]{d62728}{\textbf{4.2\%}}\\
\bottomrule
\end{longtabu}
\endgroup{}

\textbf{Top 3 rows (highlighted)} all share 110 km/h + Dusk/Darkness +
Country non-urban. Fatal rates of 6.7-7.8\% are 5x the overall baseline
of \textasciitilde1.4\%.

\textbf{Feature validation:} All three risk dimensions - speed
(\texttt{speed\_band}), lighting (\texttt{is\_night}), and location
(\texttt{Urbanisation}) - are directly captured in our engineered
feature set.

\textbf{Interaction matters:} 100 km/h + Daylight + Rural still yields
5.5\% Fatal. Speed alone is not sufficient - the model needs to learn
joint combinations.
\end{frame}

\begin{frame}[fragile]{EDA \& Data Understanding: Temporal Patterns}
\phantomsection\label{eda-data-understanding-temporal-patterns}
\pandocbounded{\includegraphics[keepaspectratio]{STAT5003_W12G07_A2_Presentation_files/figure-beamer/slide-temporal-1.pdf}}

\begin{quote}
\textbf{Key finding:} Volume peaks at weekday commute hours, motivating
\texttt{is\_rush\_hour} and \texttt{is\_weekend} features. Fatal
\emph{rate} (not count) is elevated on weekends and late-night hours.
\end{quote}
\end{frame}

\begin{frame}[fragile]{Data Preprocessing: Pipeline}
\phantomsection\label{data-preprocessing-pipeline}
\textbf{Steps 1-5: Cleaning \& Feature Engineering}

\begin{enumerate}
\tightlist
\item
  \textbf{TU Integration} - aggregated traffic-unit flags per crash
  (\texttt{has\_motorcycle}, \texttt{has\_heavy\_vehicle},
  \texttt{has\_pedestrian}, \texttt{has\_cyclist}) using
  \texttt{group\_by\ +\ summarise\ +\ left\_join}
\item
  \textbf{Feature removal} - dropped leakage columns (post-crash injury
  counts), identifiers, and variables with \textgreater75\% missing
  values; dropped \texttt{Degree\ of\ crash\ -\ detailed} (redundant
  with target)
\item
  \textbf{Imputation} - removed rows \textgreater50\% missing;
  categorical : \texttt{"Unknown"}; numeric : median
\item
  \textbf{Feature engineering} - \texttt{is\_weekend},
  \texttt{is\_rush\_hour}, \texttt{is\_night}, \texttt{is\_wet}, speed
  bands - all motivated by EDA findings
\item
  \textbf{Outlier handling} - NSW geographic bounds filter; winsorise
  \texttt{Distance} at 99th percentile
\end{enumerate}

\textbf{Steps 6-9: Modelling Setup}

\textbf{Encoding} - trees use native R \texttt{factor}; KNN/LR use
one-hot; rare LGA levels : \texttt{"Other"}

\textbf{Scaling} - Z-score standardisation within fold for KNN; trees
are scale-invariant

\textbf{Class imbalance} - inverse-frequency class weights (all models);
SMOTE applied \textbf{inside training fold only}

\textbf{Split} - 2024 data sealed as temporal holdout; 80/20 stratified
split on 2020-2023; 5-fold CV
\end{frame}

\begin{frame}{Data Preprocessing: Feature Engineering \& Class
Balancing}
\phantomsection\label{data-preprocessing-feature-engineering-class-balancing}
Engineered features - all motivated by EDA

Feature

Derivation

EDA motivation

speed\_band

Cut Speed limit: ≤50 / 60-80 / 90-100 / 110

Fatal rate rises sharply above 100 km/h

is\_night

Natural lighting in \{Darkness, Dawn, Dusk\}

Dark conditions double fatal risk

is\_rush\_hour

Two-hour intervals in \{7-9, 16-18\}

Commute peaks dominate crash volume

is\_weekend

Day of week in Saturday, Sunday

Weekend crashes differ in severity profile

is\_wet

Surface condition contains ``Wet''

Wet roads linked to skidding fatalities

has\_motorcycle

TU-level aggregated flag

Motorcyclists over-represented in fatals

has\_pedestrian

TU-level aggregated flag

Pedestrian crashes nearly always severe

Handling class imbalance

Issue

Strategy

Fatal only \textasciitilde1.4\% of crashes

Naive accuracy-maximising models predict majority class only - useless
for our goal

Class weights (all models)

Inverse-frequency weights penalise Fatal misclassification more heavily
during training

SMOTE (LR, KNN)

Synthetic oversampling applied inside training fold only; Fatal raised
to \textasciitilde10\% - test set unmodified

Primary metrics

Macro-F1 (equal weight across classes) + Fatal Recall (hard constraint:
target ≥ 0.50)
\end{frame}

\begin{frame}[fragile]{Model Fitting: Approach \& Setup}
\phantomsection\label{model-fitting-approach-setup}
We implemented six diverse classification models to balance
interpretability and predictive power. All models were trained on the
same stratified 50\% subsample (\(n \approx 55,000\)) to ensure fair
comparison.

\begin{longtable}[]{@{}
  >{\raggedright\arraybackslash}p{(\linewidth - 6\tabcolsep) * \real{0.2500}}
  >{\raggedright\arraybackslash}p{(\linewidth - 6\tabcolsep) * \real{0.2500}}
  >{\raggedright\arraybackslash}p{(\linewidth - 6\tabcolsep) * \real{0.2500}}
  >{\raggedright\arraybackslash}p{(\linewidth - 6\tabcolsep) * \real{0.2500}}@{}}
\toprule\noalign{}
\begin{minipage}[b]{\linewidth}\raggedright
Model
\end{minipage} & \begin{minipage}[b]{\linewidth}\raggedright
R Package
\end{minipage} & \begin{minipage}[b]{\linewidth}\raggedright
Imbalance Strategy
\end{minipage} & \begin{minipage}[b]{\linewidth}\raggedright
Primary Hyperparameter
\end{minipage} \\
\midrule\noalign{}
\endhead
\textbf{Multinomial LR} & \texttt{nnet} & SMOTE (inside CV) &
\texttt{decay} (L2 Penalty) \\
\textbf{Decision Tree} & \texttt{rpart} & Class Weights & \texttt{cp}
(Complexity) \\
\textbf{Random Forest} & \texttt{randomForest} & Class Weights &
\texttt{mtry} \\
\textbf{XGBoost} & \texttt{xgboost} & SMOTE / Weights &
\texttt{nrounds}, \texttt{max\_depth} \\
\textbf{KNN} & \texttt{kknn} & SMOTE (inside CV) & \texttt{k}
(Neighbors) \\
\textbf{Naive Bayes} & \texttt{naive\_bayes} & SMOTE (inside CV) &
\texttt{laplace}, \texttt{usekernel} \\
\bottomrule\noalign{}
\end{longtable}

\textbf{Validation Strategy:} - {[}cite\_start{]}\textbf{Primary
Metric:} Macro-F1 (robust to class imbalance){[}cite: 14{]}. -
{[}cite\_start{]}\textbf{Constraint:} Fatal Recall \(\geq\) 0.50
(safety-first approach){[}cite: 50{]}. -
{[}cite\_start{]}\textbf{Method:} 5-fold Cross-Validation (OOB for
Random Forest){[}cite: 49{]}.

\emph{All models trained via \texttt{caret} with 5-fold CV. Primary
metric: \textbf{Macro-F1}. Hard constraint: \textbf{Fatal Recall ≥
0.50}.}
\end{frame}

\begin{frame}{Model Fitting: Logistic Regression}
\phantomsection\label{model-fitting-logistic-regression}
Logistic regression served as our \textbf{interpretable baseline} model.
It helps show how key predictors are associated with crash severity,
while also giving us a benchmark for comparing more flexible models.

Metric

Value

Macro-F1

0.4446

Fatal Recall

0.7660

Weighted F1

0.5987

Kappa

0.2821

Key message: High coefficients for Pedestrians, Night-time, and
Non-urban areas prove our baseline LR successfully prioritises the
minority class.
\end{frame}

\begin{frame}{Model Fitting: Decision Tree}
\phantomsection\label{model-fitting-decision-tree}
Decision Tree was included because it produces clear \textbf{if-then
style rules}, making it easier to explain the model to a non-technical
audience. It reveals hierarchical relationships between risk factors.

Metric

Value

Macro-F1

0.4005

Fatal Recall

0.7660

Weighted F1

0.5349

Kappa

0.2092

Key message: The tree identifies Street Lighting and Distance as the
primary root splits. Using heavy class weights ensures the model
actively searches for Fatal outcomes, satisfying safety constraints.
\end{frame}

\begin{frame}{Model Fitting: Random Forest}
\phantomsection\label{model-fitting-random-forest}
Random Forest was chosen as our \textbf{ensemble learner} - by averaging
hundreds of decorrelated trees it reduces variance and handles the
non-linear interactions between speed, location, and lighting that
single trees miss. Built-in OOB error replaces CV, making it feasible on
large data.

Top predictors: Street lighting, First impact type, Other TU type,
Distance, Key TU type - structural \& environmental factors dominate.

Metric

Value

OOB Error Rate

22.85\%

Macro-F1

0.7029

Fatal Recall

0.4681

Weighted F1

0.7802

Kappa

0.5436
\end{frame}

\begin{frame}[fragile]{Model Fitting: XGBoost}
\phantomsection\label{model-fitting-xgboost}
XGBoost was chosen as our \textbf{boosting-based ensemble learner} - it
builds trees sequentially, where each new tree focuses on correcting the
errors made by previous trees. This makes it effective for capturing
complex non-linear interactions between crash conditions such as speed,
lighting, road surface, location, and traffic unit type. Due to the
highly imbalanced Fatal class, the model was fitted directly using
\texttt{xgboost::xgb.train()} and evaluated on the held-out test set.

Top predictors: The most important variables by Gain indicate which
crash, location, time, and road-condition features contributed most to
reducing classification error in the boosted trees.

Metric

Value

Accuracy

0.332

Macro-F1

0.266

Fatal Recall

0.33

Weighted F1

0.399

Kappa

0.003
\end{frame}

\begin{frame}{Model Fitting: KNN}
\phantomsection\label{model-fitting-knn}
\textbf{Model Overview}

\begin{itemize}
\item
  Instance-based: predicts by majority vote of k nearest neighbours
\item
  No training phase --- all computation at prediction time
\item
  Features Z-score normalised within each CV fold
\item
  Baseline to test: \emph{do similar crash conditions share severity?}
\end{itemize}

Metric

Value

Macro-F1

0.5304

Fatal Recall

0.5904

Weighted F1

0.6734

Kappa

0.3400

Note: k = round(sqrt(p)) = 18 (p = 325 one-hot features). Fatal Recall
(0.59) clears the 0.50 safety threshold, but Macro-F1 trails Random
Forest because Euclidean distance treats one-hot categoricals poorly.
\end{frame}

\begin{frame}{Model Fitting: Naive Bayes}
\phantomsection\label{model-fitting-naive-bayes}
\textbf{Why Naive Bayes?}

\begin{itemize}
\item
  Probabilistic baseline model
\item
  Suitable for categorical predictors
\item
  Fast to train on large data
\item
  Strong sensitivity to Fatal crashes
\end{itemize}

Metric

Value

Macro-F1

0.243

Fatal Recall

0.771

Weighted F1

0.249

Kappa

0.112

High Fatal Recall, but many false alarms.
\end{frame}

\begin{frame}[fragile]{Model Improvement: Hyperparameter Tuning}
\phantomsection\label{model-improvement-hyperparameter-tuning}
\begingroup\fontsize{12}{14}\selectfont

\begin{longtabu} to \linewidth {>{}l>{\raggedright}X>{\raggedright}X>{\raggedright}X>{\raggedright}X>{\raggedright}X}
\toprule
\cellcolor[HTML]{2E74B5}{\textcolor{white}{\textbf{Model}}} & \cellcolor[HTML]{2E74B5}{\textcolor{white}{\textbf{Parameter}}} & \cellcolor[HTML]{2E74B5}{\textcolor{white}{\textbf{Search Range}}} & \cellcolor[HTML]{2E74B5}{\textcolor{white}{\textbf{Best Value}}} & \cellcolor[HTML]{2E74B5}{\textcolor{white}{\textbf{Macro-F1: Before → After}}} & \cellcolor[HTML]{2E74B5}{\textcolor{white}{\textbf{Fatal Recall: Before → After}}}\\
\midrule
\textbf{Logistic Regression} & decay & 0 to 0.1 & decay = 0.001 & 0.4446 → 0.4442 & 0.7660 → 0.7553\\
\textbf{Decision Tree} & cp (complexity) & xerror minimisation & cp = 0.001 & 0.4005 → 0.4369 & 0.7660 → 0.8138\\
\textbf{\cellcolor[HTML]{eaf3fb}{Random Forest}} & \cellcolor[HTML]{eaf3fb}{mtry} & \cellcolor[HTML]{eaf3fb}{tuneRF: 4, 6, 9, 13} & \cellcolor[HTML]{eaf3fb}{mtry = 13} & \cellcolor[HTML]{eaf3fb}{0.7029 → 0.7079} & \cellcolor[HTML]{eaf3fb}{0.4681 → 0.5106}\\
\textbf{XGBoost} & nrounds / max\_depth / eta & 100-500 / 3-6 / 0.03-0.3 & best grid point & 0.266 → 0.267 & 0.33 → 0.34\\
\textbf{KNN} & k (neighbours) & k = round(sqrt(p)) & k = 18 & 0.5304 → 0.5304 & 0.5904 → 0.5904\\
\addlinespace
\textbf{Naive Bayes} & usekernel / laplace & kernel: F / laplace: 0, 0.5, 1 & laplace = 0 & 0.243 → 0.245 & 0.771 → 0.782\\
\bottomrule
\end{longtabu}
\endgroup{}

\begin{quote}
RF: \texttt{tuneRF} identified \textbf{mtry = 13} (default = 6) - Fatal
Recall climbed from \textbf{0.4681 → 0.5106}, crossing the 0.50 safety
constraint (Macro-F1 +0.005). DT: pruning at \textbf{cp = 0.001} lifted
Fatal Recall from \textbf{0.7660 → 0.8138} with Macro-F1 +0.036.
\end{quote}

\textbf{Where tuning paid off:} DT and RF gained the most. Pruning at
\texttt{cp\ =\ 0.001} let the tree grow deep enough to isolate Fatal
regions, lifting \textbf{both} Macro-F1 (+0.036) \textbf{and} Fatal
Recall (+0.048). RF's larger \texttt{mtry\ =\ 13} forced each split to
consider more candidate variables, breaking the dominance of
Injury-associated features and pushing Fatal Recall across the 0.50
safety threshold.

\textbf{Where tuning barely moved the needle:} LR (decay), XGBoost
(depth/eta/rounds), and NB (laplace) each shifted Macro-F1 by less than
±0.005. This suggests the baseline configurations already sit close to
each model's capacity on this feature set - further gains require
\textbf{structural changes} (e.g., richer encoding, calibrated
thresholds, different model class), not more tuning.

\textbf{Methodology note:} All tuning was driven by cross-validated
\textbf{Macro-F1} (or logLoss for NB), never by Fatal Recall directly.
Fatal Recall ≥ 0.50 is treated as a \textbf{hard post-hoc constraint}
rather than an optimisation target - this avoids inflating Fatal Recall
at the cost of collapsing precision on the other two classes.
\end{frame}

\begin{frame}{Model Evaluation}
\phantomsection\label{model-evaluation}
\begingroup\fontsize{13}{15}\selectfont

\begin{longtabu} to \linewidth {>{}l>{\raggedright}X>{}l>{\raggedright}X>{\raggedright}X>{\raggedright}X}
\toprule
\cellcolor[HTML]{2E74B5}{\textcolor{white}{\textbf{Model}}} & \cellcolor[HTML]{2E74B5}{\textcolor{white}{\textbf{Macro-F1}}} & \cellcolor[HTML]{2E74B5}{\textcolor{white}{\textbf{Fatal Recall}}} & \cellcolor[HTML]{2E74B5}{\textcolor{white}{\textbf{Weighted F1}}} & \cellcolor[HTML]{2E74B5}{\textcolor{white}{\textbf{Kappa}}} & \cellcolor[HTML]{2E74B5}{\textcolor{white}{\textbf{Meets Fatal Recall >= 0.50}}}\\
\midrule
\textbf{Logistic Regression (tuned)} & 0.4442 & \textcolor[HTML]{d62728}{0.7553} & 0.5987 & 0.2821 & Yes\\
\textbf{Decision Tree (tuned)} & 0.4369 & \textcolor[HTML]{d62728}{0.8138} & 0.5854 & 0.2755 & Yes\\
\textbf{\cellcolor[HTML]{eaf3fb}{\textbf{Random Forest (tuned)}}} & \cellcolor[HTML]{eaf3fb}{\textbf{0.7079}} & \textcolor[HTML]{d62728}{\cellcolor[HTML]{eaf3fb}{\textbf{0.5106}}} & \cellcolor[HTML]{eaf3fb}{\textbf{0.7817}} & \cellcolor[HTML]{eaf3fb}{\textbf{0.5473}} & \cellcolor[HTML]{eaf3fb}{\textbf{Yes}}\\
\textbf{XGBoost (tuned)} & 0.267 & \textcolor[HTML]{d62728}{0.34} & 0.401 & 0.004 & No\\
\textbf{KNN (k = 18)} & 0.5304 & \textcolor[HTML]{d62728}{0.5904} & 0.6734 & 0.3400 & Yes\\
\addlinespace
\textbf{Naive Bayes (tuned)} & 0.245 & \textcolor[HTML]{d62728}{0.782} & 0.253 & 0.112 & Yes\\
\bottomrule
\end{longtabu}
\endgroup{}

\begin{quote}
\textbf{Random Forest (tuned)} is the only model with Macro-F1
\textgreater{} 0.55 that \emph{also} satisfies Fatal Recall
\textgreater= 0.50. LR/DT have higher Fatal Recall (\textgreater0.75)
but Macro-F1 \textless{} 0.45 - too many false alarms. XGBoost fails the
safety floor. NB matches LR's Fatal Recall but Macro-F1 \textless{}
0.25.
\end{quote}

\textbf{Eliminated by the safety floor:} XGBoost is the only model below
the hard constraint (Fatal Recall = 0.34 \textless{} 0.50). Despite
extensive grid search, its weighted-loss configuration converged to a
near-trivial classifier - Kappa ≈ 0.004 indicates predictions are barely
better than random for this multi-class task.

\textbf{Pass the floor, but Macro-F1 too low:} LR, DT, and NB all clear
Fatal Recall \textgreater{} 0.75 but stay below Macro-F1 = 0.45 - they
detect fatals aggressively at the cost of precision on Injury and
Non-casualty. Useful as high-recall safety nets, but not balanced enough
as primary classifiers.

\textbf{Pass both gates - the shortlist:} Only KNN (Macro-F1 = 0.5304)
and RF (Macro-F1 = 0.7079) clear \emph{both} thresholds. RF leads KNN by
\textbf{+0.18 Macro-F1}, with consistent gaps on Weighted F1 (+0.11) and
Kappa (+0.21) - the ranking is robust across all four metrics, not a
single-metric artefact.
\end{frame}

\begin{frame}[fragile]{Model Selection}
\phantomsection\label{model-selection}
\textbf{Winner: Random Forest (tuned, mtry = 13)}

Only model that simultaneously:

\begin{itemize}
\tightlist
\item
  Achieves the \textbf{highest Macro-F1 (0.7079)} across all six models
\item
  Achieves the \textbf{highest Weighted F1 (0.7817)} and \textbf{Kappa
  (0.5473)}
\item
  \textbf{Satisfies the hard Fatal Recall constraint (0.5106
  \textgreater= 0.50)}
\end{itemize}

NB \& LR have higher Fatal Recall but collapse on Macro-F1 (massive
false alarm rates). XGBoost fails the safety constraint. KNN meets it
but loses 0.18 Macro-F1 vs RF.

\textbf{Top predictors (RF Variable Importance):}

\begin{longtable}[]{@{}ll@{}}
\toprule\noalign{}
Rank & Feature \\
\midrule\noalign{}
\endhead
1 & Street lighting \\
2 & First impact type \\
3 & Other TU type \\
4 & Distance \\
5 & Key TU type \\
6 & Speed limit / \texttt{speed\_band} \\
7 & \texttt{has\_pedestrian} \\
\bottomrule\noalign{}
\end{longtable}

Structural \& environmental factors dominate over temporal ones -
aligning with the EDA finding that 110 km/h + Darkness + Country
non-urban scenarios drive \textasciitilde5x baseline Fatal risk.
\end{frame}

\begin{frame}{Final Results \& Conclusion}
\phantomsection\label{final-results-conclusion}
\textbf{Research Question - Answered:}

\begin{quote}
\textbf{Yes} - environmental and traffic-unit features can predict
3-class crash severity. Tuned Random Forest reaches Macro-F1 = 0.7079
and correctly flags 51\% of Fatal crashes on the held-out test set, with
an OOB error rate of 22.42\%.
\end{quote}

\textbf{Key Findings:}

\begin{itemize}
\tightlist
\item
  \textbf{Speed band, street lighting, and impact type} are the
  strongest predictors - mirroring EDA
\item
  Class imbalance handling matters more than algorithm choice:
  undersampling + class weights lifted Fatal Recall from
  \textasciitilde0 (naive) to 0.51 on RF
\item
  \textbf{Pedestrian and motorcycle involvement} are top binary flags
  for fatal outcomes
\item
  Tree ensembles dominate parametric and instance-based models on this
  problem
\end{itemize}

\textbf{Practical implications:}

\begin{longtable}[]{@{}
  >{\raggedright\arraybackslash}p{(\linewidth - 2\tabcolsep) * \real{0.5000}}
  >{\raggedright\arraybackslash}p{(\linewidth - 2\tabcolsep) * \real{0.5000}}@{}}
\toprule\noalign{}
\begin{minipage}[b]{\linewidth}\raggedright
Stakeholder
\end{minipage} & \begin{minipage}[b]{\linewidth}\raggedright
Action
\end{minipage} \\
\midrule\noalign{}
\endhead
\textbf{Traffic Authorities} & Target lighting upgrades + speed reviews
on 90-110 km/h rural corridors flagged by the model \\
\textbf{Emergency Responders} & Use the model's Fatal probability score
to pre-allocate trauma resources before scene arrival \\
\textbf{Policymakers} & Evidence for stricter rural-night-speed policy;
quantified marginal risk of pedestrian + heavy-vehicle interactions \\
\textbf{Insurers} & Inputs for granular, condition-aware risk pricing \\
\bottomrule\noalign{}
\end{longtable}

The 0.7079 Macro-F1 is the highest among the six candidates while still
meeting the safety-first 0.50 Fatal Recall floor - a balance no other
model achieved.
\end{frame}

\begin{frame}{Reflection \& Future Directions}
\phantomsection\label{reflection-future-directions}
\textbf{Limitations}

\begin{itemize}
\tightlist
\item
  \textbf{Random undersampling discards \textasciitilde90,000
  majority-class rows} - introduces variance and a distribution shift
  relative to the true operating environment
\item
  \textbf{XGBoost underperformed unexpectedly} (Macro-F1
  \textasciitilde0.27) - likely due to one-hot expansion of
  high-cardinality categoricals interacting poorly with the chosen
  weight scheme; needs revisiting
\item
  \textbf{Naive Bayes' conditional independence assumption is violated}
  (e.g., speed and urbanisation are highly correlated)
\item
  \textbf{No temporal holdout in final report run} - 2024 data prepared
  but not yet evaluated; potential temporal drift unmeasured
\item
  Free-text fields (Street, Town, Identifying feature) were dropped -
  could carry signal
\end{itemize}

\textbf{Future Directions}

\begin{itemize}
\tightlist
\item
  \textbf{Wider mtry / ntree search} for RF - OOB curve had not
  converged at ntree = 100, so further gains likely
\item
  \textbf{Target encoding} for high-cardinality categoricals (LGA, TU
  type) instead of one-hot - especially for XGBoost
\item
  \textbf{Cost-sensitive learning} with explicit asymmetric
  misclassification costs (cost of missing a Fatal
  \textgreater\textgreater{} cost of false alarm)
\item
  \textbf{Calibrated probabilities + threshold tuning} rather than
  argmax classification - lets stakeholders pick their own operating
  point
\item
  \textbf{Geospatial features} from Lat/Lon (proximity to intersections,
  road class, road geometry)
\item
  \textbf{Evaluate on 2024 temporal holdout} to quantify drift before
  deployment
\end{itemize}
\end{frame}

\end{document}
